\documentclass[a4j,  9pt]{ltjsarticle}

\usepackage{../mypackage}

\begin{document}

  \section{写像}
    前節では集合論の概要を説明した。ついでここでは集合間を関係づける概念, 写像についてみておこう。
    
    \subsection{写像の定義}
        $\ds X,  Y$を空でない集合とする。\par
        任意の$\ds x \in X$に対し, ある$\ds y \in Y$が一意に定まるとき, このような\textbf{対応規則}を\textbf{$\ds X$から$\ds Y$への写像}と呼ぶ。\par
        また\textbf{$\ds X,  Y$が数を要素とする集合}の場合特に\textbf{関数}と呼ぶ。つまり$\ds 関数 \subset 写像$である。\par
        簡潔に表現すれば以下。

        \columnbreak

        \begin{equation*}
          \begin{split}
            $\ds f:X \mapsto Y \defineProposition \forall_{x \in X} [\exists!_{y \in Y} [y = f(x)]]$
          \end{split}
        \end{equation*}

        \begin{equation*}
          \begin{tikzpicture}[auto]
          
            %Xのベン図
            \fill[color=black!50!white,  draw=black] (0, 0) ellipse (20pt and 40pt);
            \node[set label, text=black] (X) at (0,  40pt) {$\ds X$};

            %Yのベン図
            \draw (80pt, 0) ellipse (20pt and 40pt);
            \fill[color=black!50!white,  draw=black] (80pt,  0) ellipse (10pt and 20pt); % black!50!white で黒と白50%
            \node[set label, text=black] (Y) at (80pt,  40pt) {$\ds Y$};
            \node[set label, text=black] (Y') at (80pt,  20pt) {$\ds Y'$};

            %写像f
            \draw[->] (0,  0) to node {$\scriptstyle f$} (80pt,  0);

          \end{tikzpicture}
        \end{equation*}

      \subsection{定義域・値域}
        $\ds f:X \mapsto Y$において

        \begin{cases}
          $\ds X$を\textbf{$\ds f$の定義域}\\
          $\ds Y'(\subset Y)$を\textbf{$\ds f$の値域}\\
        \end{cases}

        という。\par

        \columnbreak

        値域$\ds R_f$の具体的な定義は以下。
      
        \columnbreak

        \begin{equation*}
          \begin{split}
            $\ds R_f \defineProposition \{y \in Y \mid \exists_x [x \in X \land y=f(x)]\}$
          \end{split}
        \end{equation*}

        % \footnote{
        %   $\ds \exists_x [x \in X \land y=f(x)]$は自由変項$\ds y$が残るため$\ds y$の条件と言える。
        % }

      \subsection{単射（行先がかぶらない写像）}
        写像にはいくつか特性が存在する。その特性の一つを有する写像が\textbf{単射}と呼ばれるものである。\par
        \textbf{$\ds f:X \mapsto Y$が単射である}とは, \textbf{任意の$x, x'\in X$に対して「$\ds x \ne x' \Longrightarrow f(x) \ne f(x')$」が成立すること}である。\par
        つまり簡潔に表せば以下。

        \begin{equation*}
          \begin{split}\notag
            $\ds f:X \mapsto Yが単射$     & $\defineProposition      \forall_{x, x'\in X} [x \ne x' \Longrightarrow f(x) \ne f(x')]$\\
                                          & $\ds \Longleftrightarrow  \forall_{x, x'\in X} [f(x) = f(x') \Longrightarrow x = x']$\\
          \end{split}
        \end{equation*}

        \begin{equation*}
          \begin{tikzpicture}[auto]
          
            %Xのベン図
            \fill[color=black!50!white,  draw=black] (0, 0) ellipse (20pt and 40pt);
            \node[set label,  text=black] (X) at (0,  40pt) {$\ds X$};

            %Yのベン図
            \draw (80pt, 0) ellipse (20pt and 40pt);
            \fill[color=black!50!white,  draw=black] (80pt,  0) ellipse (10pt and 20pt); % black!50!white で黒と白50%
            \node[set label,  text=black] (Y) at (80pt,  40pt) {$\ds Y$};

            %x->y
            \node (x) at (0,  15pt) {$\ds x$};
            \node (y) at (80pt,  15pt) {$\ds y$};
            \draw[->] (0,  10pt) to node {$\scriptstyle f(x)$} (80pt,  10pt);

            %x'->y'
            \node (x') at (0,  -3pt) {$\ds x'$};
            \node (y') at (80pt,  -3pt) {$\ds y'$};
            \draw[->] (0,  -10pt) to node {$\scriptstyle f(x')$} (80pt,  -10pt);

          \end{tikzpicture}
        \end{equation*}

      \subsection{全射}
        写像の特使の一つを有するものを単射と言った。他にも\textbf{全射}と呼ばれるものもある。\par
        \textbf{$\ds f:X \mapsto Y$が全射である}とは, \textbf{任意の$y \in Y$に対して, ある$\ds x \in X$が存在し$\ds f(x)=y$となること}である。\par
        つまり簡潔に表せば以下。

        \begin{equation*}
          \begin{split}\notag
            $\ds f:X \mapsto Yが全射$     & $\defineProposition      \forall_{y\in Y} [\exists_{x\in X} [f(x) = y]]$\\
          \end{split}
        \end{equation*}

        \begin{center}
          \begin{tikzpicture}[auto]
          
            %Xのベン図
            \draw (0, 0) ellipse (20pt and 40pt);
            \fill[color=black!50!white,  draw=black] (0,  0) ellipse (10pt and 20pt); % black!50!white で黒と白50%
            \node[set label,  text=black] (X) at (0,  40pt) {$\ds X$};

            %Yのベン図
            \fill[color=black!50!white,  draw=black] (80pt, 0) ellipse (20pt and 40pt);
            \node[set label,  text=black] (Y) at (80pt,  40pt) {$\ds Y$};

            %x->y
            \node (x) at (0,  15pt) {$\ds x$};
            \node (y) at (80pt,  15pt) {$\ds y$};

            %x'->y'
            \node (x') at (0,  5pt) {$\ds x'$};
            \node (y') at (80pt,  -3pt) {$\ds y'$};

            %x''->y'
            \node (x'') at (0,  -13pt) {$\ds x''$};

            \draw[->] (0,  10pt) to node {$\scriptstyle f(x)$} (80pt,  10pt);
            \draw[->] (0,  0) to node {$\scriptstyle f(x')$} (80pt,  -10pt);
            \draw[->] (0,  -20pt) to node {$\scriptstyle f(x'')$} (80pt,  -10pt);

          \end{tikzpicture}
        \end{center}

        \subsection{全単射（1対1対応, 変換）}
        既に写像の特性を持つものを2つ紹介したが, これらの特性を同時に有するものも存在する。それを\textbf{全単射}といい, \textbf{単射と全射それぞれの定義から成る合成命題}で定義される。\par
        つまり

        \begin{equation*}
          \begin{split}
            \textbf{$\ds f:X \mapsto Y$が全単射である}とは, \textbf{$\ds f$が単射かつ全射であるとき}をいう。
          \end{split}
        \end{equation*}
        
        また, \textbf{定義域と値域が等しい}ことよりこの対応を\textbf{1対1対応}とも呼び, さらに\textbf{互いの集合間をすべての要素が行き来できる}ことより\textbf{変換}とも呼ぶ。(ただし変換に関しては必ずしも1対1対応でなくてもよい。)\par
        つまり簡潔に表せば以下。

        \begin{equation*}
          \begin{split}
            $\ds f:X \mapsto Yが全単射 \defineProposition \forall_{x, x'\in X} [f(x) = f(x') \Longrightarrow x = x'] \land \forall_{y\in Y} [\exists_{x\in X} [f(x) = y]]$
          \end{split}
        \end{equation*}

        \begin{equation*}
          \begin{tikzpicture}[auto]
          
            %Xのベン図
            \fill[color=black!50!white,  draw=black] (0, 0) ellipse (20pt and 40pt);
            \node[set label,  text=black] (X) at (0,  40pt) {$\ds X$};

            %Yのベン図
            \fill[color=black!50!white,  draw=black] (80pt, 0) ellipse (20pt and 40pt);
            \node[set label,  text=black] (Y) at (80pt,  40pt) {$\ds Y$};

            %x->y
            \node (x) at (0,  15pt) {$\ds x$};
            \node (y) at (80pt,  15pt) {$\ds y$};
            \draw[->] (0,  10pt) to node {$\scriptstyle f(x)$} (80pt,  10pt);

            %x'->y'
            \node (x') at (0,  -3pt) {$\ds x'$};
            \node (y') at (80pt,  -3pt) {$\ds y'$};
            \draw[->] (0,  -10pt) to node {$\scriptstyle f(x')$} (80pt,  -10pt);

          \end{tikzpicture}
        \end{equation*}
        
\end{document}